\begin{ExerciseBlock}
\ExerciseHeading{}{Verständnisfragen}{rigid_bodies_checks}
Entscheiden Sie, welche Aussagen wahr oder falsch sind, und begründen Sie Ihre Entscheidung.
\begin{enumerate}
\item Ein starrer Körper dehnt sich aus, wenn er sehr schnell rotiert wird.\TrueFalseAnswer{falsch}
\item Der Schwerpunkt eines starren Körpers mit bekannter Massendichte lässt sich durch Volumenintegration bestimmen.\TrueFalseAnswer{wahr}
\item Bei Drehung um eine feste Achse gelten für die achsparallele Komponente des Drehimpulses $L_\parallel=J\omega$ und für die Rotationsenergie $E_\mathrm{rot}=J\omega^2/2$.\TrueFalseAnswer{wahr}
\item Das Trägheitsmoment $J$ ist unabhängig von der Dichte eines Materials.\TrueFalseAnswer{falsch}
\item Ein Quader habe unterschiedlich lange Seiten $a<b<c$ und rotiere um seinen Schwerpunkt. Das kleinste Trägheitsmoment $J$ hat der Quader, wenn die Drehachse $\vec{\omega}$ parallel zur längsten Seite $c$ ist.\TrueFalseAnswer{wahr}
\item Kennt man den Schwerpunkt eines starren Körpers, so erlaubt der Steinersche Satz die Berechnung des Trägheitsmoments bei Rotation um eine beliebige Achse.\TrueFalseAnswer{falsch}
\item Die kinetische Energie eines starren Körpers setzt sich aus der Translationsenergie des Schwerpunkts und der Rotationsenergie um den Schwerpunkt zusammen.\TrueFalseAnswer{wahr}
\item Bei gleicher Masse und gleichem Rollradius beschleunigt ein Körper mit größerem Trägheitsmoment auf einer schiefen Ebene stärker.\TrueFalseAnswer{falsch}
\item Bei gleicher Masse, gleichem Rollradius und gleicher Anfangsgeschwindigkeit erreicht der Körper mit größerem Trägheitsmoment beim verlustfreien Hinaufrollen eine größere Höhe.\TrueFalseAnswer{wahr}
\end{enumerate}
\end{ExerciseBlock}
